% manuscript-sha256: 2a8bfe8f3374c09da1b75356a37c983f3222d69b4bb1a8b23046cd5957793d38
% Options for packages loaded elsewhere
\PassOptionsToPackage{unicode}{hyperref}
\PassOptionsToPackage{hyphens}{url}
\PassOptionsToPackage{dvipsnames,svgnames,x11names}{xcolor}
\documentclass[
  11pt,
]{article}
\usepackage{xcolor}
\usepackage[margin=25mm]{geometry}
\usepackage{amsmath,amssymb}
\setcounter{secnumdepth}{-\maxdimen} % remove section numbering
\usepackage{iftex}
\ifPDFTeX
  \usepackage[T1]{fontenc}
  \usepackage[utf8]{inputenc}
  \usepackage{textcomp} % provide euro and other symbols
\else % if luatex or xetex
  \usepackage{unicode-math} % this also loads fontspec
  \defaultfontfeatures{Scale=MatchLowercase}
  \defaultfontfeatures[\rmfamily]{Ligatures=TeX,Scale=1}
\fi
\usepackage{lmodern}
\ifPDFTeX\else
  % xetex/luatex font selection
  \setmainfont[]{Latin Modern Roman}
  \setmonofont[]{Latin Modern Mono}
\fi
% Use upquote if available, for straight quotes in verbatim environments
\IfFileExists{upquote.sty}{\usepackage{upquote}}{}
\IfFileExists{microtype.sty}{% use microtype if available
  \usepackage[]{microtype}
  \UseMicrotypeSet[protrusion]{basicmath} % disable protrusion for tt fonts
}{}
\makeatletter
\@ifundefined{KOMAClassName}{% if non-KOMA class
  \IfFileExists{parskip.sty}{%
    \usepackage{parskip}
  }{% else
    \setlength{\parindent}{0pt}
    \setlength{\parskip}{6pt plus 2pt minus 1pt}}
}{% if KOMA class
  \KOMAoptions{parskip=half}}
\makeatother
\usepackage{longtable,booktabs,array}
\usepackage{calc} % for calculating minipage widths
% Correct order of tables after \paragraph or \subparagraph
\usepackage{etoolbox}
\makeatletter
\patchcmd\longtable{\par}{\if@noskipsec\mbox{}\fi\par}{}{}
\makeatother
% Allow footnotes in longtable head/foot
\IfFileExists{footnotehyper.sty}{\usepackage{footnotehyper}}{\usepackage{footnote}}
\makesavenoteenv{longtable}
\setlength{\emergencystretch}{3em} % prevent overfull lines
\providecommand{\tightlist}{%
  \setlength{\itemsep}{0pt}\setlength{\parskip}{0pt}}
\usepackage{bookmark}
\IfFileExists{xurl.sty}{\usepackage{xurl}}{} % add URL line breaks if available
\urlstyle{same}
\hypersetup{
  colorlinks=true,
  linkcolor={Maroon},
  filecolor={Maroon},
  citecolor={Blue},
  urlcolor={Blue},
  pdfcreator={LaTeX via pandoc}}

\author{}
\date{}

\usepackage{xurl}
\usepackage{newunicodechar}
\usepackage{needspace}
\newunicodechar{∎}{\ensuremath{\square}}
\newunicodechar{→}{\ensuremath{\to}}
\newunicodechar{←}{\ensuremath{\leftarrow}}
\newunicodechar{↦}{\ensuremath{\mapsto}}
\newunicodechar{≠}{\ensuremath{\ne}}
\newunicodechar{⊆}{\ensuremath{\subseteq}}
\newunicodechar{∩}{\ensuremath{\cap}}
\newunicodechar{∪}{\ensuremath{\cup}}
\newunicodechar{∅}{\ensuremath{\varnothing}}
\newunicodechar{∈}{\ensuremath{\in}}
\newunicodechar{∉}{\ensuremath{\notin}}
\newunicodechar{⊔}{\ensuremath{\sqcup}}
\newunicodechar{−}{\ensuremath{-}}
\newunicodechar{²}{\ensuremath{{}^2}}
\newunicodechar{³}{\ensuremath{{}^3}}
\newunicodechar{φ}{\ensuremath{\varphi}}
\newunicodechar{Σ}{\ensuremath{\Sigma}}
\newunicodechar{Γ}{\ensuremath{\Gamma}}
\setlength{\emergencystretch}{3em}
\begin{document}

\begin{center}
{\LARGE\bfseries From Admitted Distinctions to S(n) and MSS-N: Presentations, Retractions, and Obstruction\par}
Jeremy H. Carroll\par
2026-10-09\par
Version 1.0.0\par
\end{center}

\subsection{Abstract}\label{abstract}

An admitted distinction raises two separate questions: how should it be
represented, and which prior information must remain recoverable as the
representation changes? This paper consolidates the S(n)
presentation-tower account, the MSS-N mathematical monograph and the
corrected Pre-Genesis-through-MSS-N architectural walkthrough. It
distinguishes native admission stages, alphabet inventories, grammar
profiles and atlas targets. It develops finite demand accounting,
observation refinement, typed recovery on an old image, and conditional
claims about minimal extensions. At the language level, alphabet
inclusion admits an exact restriction retraction. At the chain-complex
level, a stronger preservation promise has a sharp limit: a chain-map
retraction makes the induced homology map injective, so an extension
killing an old nonzero class cannot retain that retraction. Explicit
graph examples show why syntactic recovery, observational discrimination
and topological repair need different checks. These distinctions provide
a bridge of questions and typed obligations into WDP; they do not claim
that the complete WDP calculus or a physical theory has been forced from
the initial inscription.

\subsection{1. What ``Before the First Mark''
supplies}\label{what-before-the-first-mark-supplies}

The primitive monograph \emph{Before the First Mark} begins before
interpretation of \texttt{\textless{}\textgreater{}}. It distinguishes
absent input, the empty word, the literal inscription and an admitted
ground. It explains why a task cannot distinguish two objects through a
report that identifies them and why a locally valid reply can still
depend on the wrong historical state.

Here those lessons become requirements on representations. A candidate
extension must name what it changes, what it preserves and which problem
it resolves. A representation's convenience, mathematical richness or
physical analogy does not authorize its admission to an earlier regime.

The consolidation adopts two architectural rules:

\begin{enumerate}
\def\labelenumi{\arabic{enumi}.}
\tightlist
\item
  A proposed native extension must be justified against an unresolved,
  declared obligation and the allowed response class.
\item
  A mathematical realization has no generative authority over that
  native regime.
\end{enumerate}

These are policies for the programme. Whether a particular
implementation enforces them requires evidence about that
implementation. A theorem about one model does not establish estate-wide
enforcement.

\subsection{2. The indices and layers}\label{the-indices-and-layers}

\begin{longtable}[]{@{}
  >{\raggedright\arraybackslash}p{(\linewidth - 4\tabcolsep) * \real{0.3333}}
  >{\raggedright\arraybackslash}p{(\linewidth - 4\tabcolsep) * \real{0.3333}}
  >{\raggedright\arraybackslash}p{(\linewidth - 4\tabcolsep) * \real{0.3333}}@{}}
\toprule\noalign{}
\begin{minipage}[b]{\linewidth}\raggedright
Symbol or object
\end{minipage} & \begin{minipage}[b]{\linewidth}\raggedright
Meaning
\end{minipage} & \begin{minipage}[b]{\linewidth}\raggedright
Separate quantity
\end{minipage} \\
\midrule\noalign{}
\endhead
\bottomrule\noalign{}
\endlastfoot
S(n) & A presentation stage in the specified admission account. & The
number of glyphs or a homotopy degree. \\
k & Alphabet size under a declared token policy. & Number of decoded
objects or operations. \\
p & Grammar/profile and public start. & Semantic interpretation or
evaluation verdict. \\
MSS-m & An externally specified atlas target. & A native Genesis stage
or a completed lift. \\
H\_j & Homology in degree j of a declared complex. & Stage n or
active-demand count. \\
C(m), D & A declared coverage measure and length of a contiguous
licensed chain. & One another; neither is inferred from the number of
folders. \\
\end{longtable}

For example, S(2) in the teaching crosswalk records retained but
observationally collapsed data. MSS-2 is a target about relations.
Equating them by their numerals silently changes the object under
discussion.

Four responsibilities remain distinct: native admissions; mathematical
representations; downstream interpretations; and verification/evidence
machinery. A Lean proof, a Python calculation and a SQL view may refer
to related structures, but shared terminology does not prove that they
evaluate the same property.

\subsection{3. The finite-stage
crosswalk}\label{the-finite-stage-crosswalk}

The
\href{../../../../../omega_engine_education/documentation/001_genesis/genesis.md}{Genesis
account} is the source for admission, observation, candidate costs and
the sealed/post-seal distinctions. The following is a presentation
crosswalk, not a claim that the displayed glyphs derive their own
semantics or that every row is forced independently of policy.

\begin{longtable}[]{@{}
  >{\raggedright\arraybackslash}p{(\linewidth - 4\tabcolsep) * \real{0.3333}}
  >{\raggedright\arraybackslash}p{(\linewidth - 4\tabcolsep) * \real{0.3333}}
  >{\raggedright\arraybackslash}p{(\linewidth - 4\tabcolsep) * \real{0.3333}}@{}}
\toprule\noalign{}
\begin{minipage}[b]{\linewidth}\raggedright
Stage
\end{minipage} & \begin{minipage}[b]{\linewidth}\raggedright
Presentation question
\end{minipage} & \begin{minipage}[b]{\linewidth}\raggedright
Required qualification
\end{minipage} \\
\midrule\noalign{}
\endhead
\bottomrule\noalign{}
\endlastfoot
S(0) & What is received as the opaque ground inscription? &
\texttt{\textless{}\textgreater{}} has no inherited object-language
meaning. Ground inertia is relative to the specified rules. \\
S(1) & Has a newly supplied mark been consumed or admitted? & Supply
alone may be an internal no-op. \\
S(2) & What data have been retained, and which profiles remain equal? &
Host-level multiplicity does not create an internal distinguishing
context. \\
S(3) & How can a supplied pair-separation task be met? & The cost-two
result uses a particular response grammar and cost; it is not a
universal witness count. \\
S(4) & Have the relevant pairs been separated under the chosen closure
policy? & Pair closure needs its finite-carrier and preservation
assumptions. \\
S(5) & Where does the next requirement or fresh datum come from? &
External continuation is a declared policy, not an endogenous generator
proved by saturation. \\
S(6) & How are old presentations represented and recovered after an
extension? & Typed maps and their preservation laws must be supplied and
checked. \\
\end{longtable}

A branching set of tied responses is not automatically a linear tower. A
chosen cost order, equivalence and continuation policy select or retain
its branches. Subsequent atlas targets may ask for graphs, algebra,
categories or geometry. Their presence in the atlas does not establish
that the native process has reached them.

\subsection{4. Active demands, receipt sets and ordered
histories}\label{active-demands-receipt-sets-and-ordered-histories}

Let W\_n be a finite register of active demands. A demand identifies its
carrier, protected condition, witness data and activity predicate. It is
not a bare label such as ``curvature'' or ``missing structure.'' Let
D\_\{n+1\} be the demands discharged in a step and B\_\{n+1\} the
demands introduced. Assume

\[
D_{n+1}\subseteq W_n,\qquad
B_{n+1}\cap(W_n\setminus D_{n+1})=\varnothing,
\]

and define

\[
W_{n+1}=(W_n\setminus D_{n+1})\sqcup B_{n+1}.
\]

\textbf{Proposition 001.1 --- finite demand accounting.} With
s\_n=\textbar W\_n\textbar, d\_\{n+1\}=\textbar D\_\{n+1\}\textbar{} and
b\_\{n+1\}=\textbar B\_\{n+1\}\textbar,

\[
s_{n+1}=s_n-d_{n+1}+b_{n+1}.
\]

\textbf{Proof.} Subset removal gives \textbar W\_n minus
D\_\{n+1\}\textbar=\textbar W\_n\textbar−\textbar D\_\{n+1\}\textbar.
Disjoint union then adds \textbar B\_\{n+1\}\textbar. ∎

This identity does not prove which demands an engine actually carries.
In particular, the previously displayed sequence (0,1,1,1,0,1,1) follows
from one chosen assignment of singleton demands to the crosswalk. It is
an illustrative ledger, not a universal native survivor sequence or a
physical observable.

A cumulative receipt set H\_\{n+1\}=H\_n union C\_\{n+1\} has
nondecreasing cardinality. An ordered event log additionally retains
order and multiplicity. Neither is W\_n. An empty active register need
not erase receipts, imply thermodynamic equilibrium, or prove that every
possible future question has been answered. Adding two map symbols does
not imply two active demands or two newly admitted carrier elements.

\subsection{5. Binary support and observation
refinement}\label{binary-support-and-observation-refinement}

An inhabited irreflexive relation R(a,b) forces a≠b, as proved in Paper
000. This gives two endpoints for that edge, not two observers or two
certificate occurrences. The pointed graph D2 with vertices 0 and 1 and
the single edge 0→1 has exactly one endpoint-preserving map to a
specified pair (a,b); that map preserves the edge. This elementary
pointed property does not make D2 initial in an unspecified category of
all possible systems.

Suppose q:A→O identifies a and b. If a supplied tag t:A→\{0,1\} has
t(a)=0 and t(b)=1, define q'(x)=(q(x),t(x)). Then q'(a)≠q'(b), while
projection to the first component recovers q. Both claims follow
directly from equality of ordered pairs. The tag and its construction
are supplied data. A method for finding tags is a separate algorithmic
obligation.

For several required pairs, form a conflict graph on A. A tag assignment
separating every edge is a proper coloring. Thus the number of tag
values needed for many pairwise constraints need not remain two: a
three-clique needs three colors. The constraints remain binary without
implying a globally binary solution space.

\subsection{6. Typed recovery on the old
image}\label{typed-recovery-on-the-old-image}

Let u:A→B be a proposed lift and let

\[
\widehat r:B\to A+\{\mathrm{NON\_DESCENT}\}
\]

be a typed recovery map. The old-image law is

\[
\widehat r(u(a))=\mathrm{inl}(a)\quad(a\in A).
\]

\textbf{Proposition 001.2 --- injectivity.} This law implies that u is
injective. If u(a)=u(a'), applying the recovery map gives
inl(a)=inl(a'), hence a=a'. ∎

The law alone says nothing about off-image inputs. A checker intended to
reject all of them needs the additional specification

\[
b\notin\operatorname{im}(u)\;\Longrightarrow\;
\widehat r(b)=\mathrm{inr}(\mathrm{NON\_DESCENT}).
\]

It also needs a way to decide the image condition within its declared
domain and budget. That is immediate for some finite presentations and
not a general computability theorem.

Two further qualifications prevent stronger promises from being
inferred:

\begin{itemize}
\tightlist
\item
  \textbf{Observation preservation needs a commuting condition.} For old
  observation q:A→O and new observation q':B→O, require q' composed with
  u=q, or provide a declared transport to the old report space.
  Recovering a term does not by itself prove that a new evaluator uses
  the old observations.
\item
  \textbf{Repair needs its own proof.} An identity embedding has a
  perfect old-image law and can leave every original obstruction
  unchanged.
\end{itemize}

The pair is not generally an involution. Moreover, u composed with r-hat
is not even a map B→B until handling of the refusal summand has been
defined. No global preservation of all certificates follows from a
successful term round trip alone.

The retained
\href{../../../../../omega_engine_education/documentation/002_oesn/lean4/SnCarrier.lean}{SnCarrier
model} does not prove this law for its entire input type. The compiled
checks \nolinkurl{refusal_token_round_trip_changes_input} and
\nolinkurl{carrier_card_table_disagrees_with_higher} in
\href{lean4/SnCarrierDischarge.lean}{SnCarrierDischarge.lean} import
that module. \texttt{sectionLift\ stage\ (.nonDescent\ ⟨()⟩)} uses the
fallback empty-body lift, and \nolinkurl{retractionDownMap} then accepts a
ground presentation instead of recovering the original refusal token.
\texttt{Carrier.higher\ 0\ 0} inhabits \texttt{Carrier\ 0}, while
\texttt{carrierCard\ 0\ =\ 0}. The table is not a cardinality theorem
about that inductive type. These compiled facts discharge the
inspection. They do not prove a stage-specific admissible domain or a
corrected cardinality statement. Those repair obligations stay open, and
the unrestricted old-image law is not claimed for this model.

\subsection{7. Formal languages and conditional expressive
capacity}\label{formal-languages-and-conditional-expressive-capacity}

A symbol system must declare terminal spellings, tokenization, grammar,
public start, layout policy, identity policy, annotation decoding and
allowed manipulations. These eight choices determine what an inscription
does.

For the familiar alphabet Gamma\_38, the glyph inventory comprises 26
lowercase Latin letters, ten decimal digits and \texttt{\textless{}},
\texttt{\textgreater{}}. Some MSS grammar profiles add punctuation or
Unicode; they are not confined to these 38 characters.

Keep three relations separate:

\[
E\subseteq V\times V,\qquad
\to_{\rm rewrite}\;\subseteq\Sigma^*\times\Sigma^*,\qquad
q:V\to O.
\]

They represent an encoded relation, rewrites between inscriptions and an
observation map respectively. A change in any one need not change the
others.

With repetition admitted, the one-letter language a+ supports positive
unary counting under the length decoder. Including epsilon gives a* and
a representation of zero. With two letters, concatenation is
noncommutative: ab and ba are different words. A declared binary decoder
can use any two distinct printed letters. A relation may be encoded by
positions or a table without an arrow glyph. These are conditional
constructions, not semantic resources supplied by alphabet size alone.

\subsubsection{7.1 Free monoids and numeral
decoders}\label{free-monoids-and-numeral-decoders}

\textbf{Theorem 001.3.} Free monoids on finite alphabets of different
sizes are not isomorphic.

\textbf{Proof.} In a free monoid, the nonidentity elements that cannot
be factored as a product of two nonidentity elements are exactly the
one-letter words. A monoid isomorphism preserves identity and this
indecomposability property in both directions. It therefore induces a
bijection of the generating alphabets. Their sizes must agree. ∎

Thus binary and decimal string monoids are not isomorphic. Decoders can
map their strings onto the same natural numbers without making the
string monoids identical. A fixed-width binary encoding of decimal
digits is injective under concatenation, but not onto the binary free
monoid.

Likewise, discrete categories with two and 26 objects are not
equivalent: all isomorphism classes are singletons, and an equivalence
would have to match those classes bijectively.

\subsubsection{7.2 Three recovery constructions with different
types}\label{three-recovery-constructions-with-different-types}

For alphabets A contained in B, language inclusion and restriction are

\[
\ell:\mathcal P(A^*)\to\mathcal P(B^*),\quad\ell(L)=L,
\qquad r(M)=M\cap A^*.
\]

Then r(ell(L))=L. This is a retraction of sets of words.

For individual words, exact recovery of old words can instead return the
same word if it is in A\emph{, and a typed refusal otherwise. A total
erasure homomorphism h:B}→A* is another choice: map each new letter to
epsilon and concatenate. It recovers every old word but loses new
letters and their positions. These are three distinct contracts, not one
overloaded ``down-map.''

No abelian category of ordinary formal languages has been specified
here. An Ext calculation cannot be inferred from language restriction.
Even in an abelian category, a split extension is one zero extension
class; it does not show that an entire Ext group vanishes.

\subsection{8. Manipulations and actual information
loss}\label{manipulations-and-actual-information-loss}

A renaming f:A→B has a two-sided inverse on the entire codomain B
exactly when it is bijective. Recovery of source letters from the image
f(A) only requires injectivity. Mapping both c and d to d collapses ac
and ad. This is independent of whether the resulting word is
grammatically well formed.

For the code a↦0 and b↦00, the different source words aa and b both
encode as 00; ab and ba both encode as 000. A decoder cannot recover the
source universally without additional retained information or a
restricted source domain.

For declared directed-edge syntax:

\begin{longtable}[]{@{}ll@{}}
\toprule\noalign{}
Inscription & Represented edge \\
\midrule\noalign{}
\endhead
\bottomrule\noalign{}
\endlastfoot
\texttt{a\textgreater{}b} & a→b \\
\texttt{b\textgreater{}a} & b→a \\
\texttt{a\textless{}b} & b→a \\
\texttt{b\textless{}a} & a→b \\
\end{longtable}

Reversing the vertex sequence and every arrow in
\texttt{a\textgreater{}b\textgreater{}c} gives
\texttt{c\textless{}b\textless{}a} and preserves its directed edge
multiset. It changes the traversal order. A checker must name whether it
protects the edge multiset, the ordered trace, annotations, or exact
source bytes. AST recovery and byte recovery are different questions.

\subsection{9. Finite graphs and the homological
obstruction}\label{finite-graphs-and-the-homological-obstruction}

For a finite directed graph, regard the underlying graph as a
one-dimensional cell complex and define

\[
C_1=\mathbb Z^E\xrightarrow{\partial_1}C_0=\mathbb Z^V,
\qquad\partial_1(e)=t(e)-s(e).
\]

If V is nonempty and the underlying undirected graph has c components,
then H\_0 is isomorphic to Z\^{}c and reduced H\_0 to Z\^{}\{c−1\}. The
maximal component count \textbar V\textbar{} occurs when no edge joins
distinct vertices; self-loops do not affect it. Adding an edge reduces
the component count only when it joins previously distinct components.
It does not separate a and b under an unchanged observation q with
q(a)=q(b).

Formal negative coefficients do not authorize reverse physical
traversal. In a CW graph, a self-loop has zero boundary and can generate
nonzero H\_1. It is not a degenerate simplex solely because its
endpoints coincide; the model matters.

\subsubsection{9.1 The chain-retraction
obstruction}\label{the-chain-retraction-obstruction}

\textbf{Theorem 001.4.} Let i:C→D and r:D→C be chain maps with r
composed with i =id\_C. Then H\_j(i) is injective for every j. In
particular, i cannot kill a nonzero homology class of C.

\textbf{Proof.} If z is a cycle in C and i(z)=partial\_D b, then

\[
z=r(i(z))=r(\partial_D b)=\partial_C r(b).
\]

So z was already a boundary in C. Therefore the kernel of H\_j(i) is
zero. ∎

This is the central preservation limit inherited from the MSS-N
monograph. It does not prohibit every useful extension; it identifies
incompatible promises. An extension may repair a cycle while preserving
inscriptions, labels or selected questions, but it cannot simultaneously
claim this chain-map retraction fixing all of C.

\subsubsection{9.2 Two isolated points}\label{two-isolated-points}

Let C\_0=Za plus Zb and C\_1=0. Add an edge e with partial\_D e=b−a,
obtaining D. A retraction fixing the original vertices would have
r\_0=id and r\_1(e)=0. The chain-map equation would require
b−a=partial\_C(0)=0, which cannot hold in the free abelian group on a
and b. The new edge kills the old reduced zero-dimensional class, so the
theorem applies.

\subsubsection{9.3 A triangle and an explicitly attached
face}\label{a-triangle-and-an-explicitly-attached-face}

Under label identity,
\texttt{a\textgreater{}b\textgreater{}c\textgreater{}a} represents edges
e\_ab, e\_bc and e\_ca. Their sum z has boundary (b−a)+(c−b)+(a−c)=0.
With no two-cells, H\_1 is generated by {[}z{]}. Adjoining a face f with
partial\_2 f=z makes z a boundary and satisfies partial\_1 partial\_2=0.
The inscription gave the one-skeleton; the face is an additional
declared object.

Killing this class prevents a chain-map retraction fixing the original
complex. It does not prove that an unrelated observational collision has
been resolved. Nor does one vanishing class show that every homology
group vanishes.

Chains and cochains must remain typed separately: C\^{}j=Hom(C\_j,Z) and
(d\^{}j phi)(c)=phi(partial\_\{j+1\}c). Filling a chain is not the same
equation as finding a cochain primitive. Homological vanishing also
differs from null-homotopy: the commutator loop in the wedge of two
circles has zero abelianized homology class but is nontrivial in its
free fundamental group.

\subsection{10. Profiles, parsers and evidence
scope}\label{profiles-parsers-and-evidence-scope}

The
\href{../../../../../omega_engine_education/documentation/003_minimal_symbol_system_n_mssn/grammar/mssn.ebnf}{MSS
grammar} defines the 22 explicit starts inventoried in Appendix A. It
supplies formation rules, not physical models or admitted semantic
operations.

\begin{longtable}[]{@{}
  >{\raggedright\arraybackslash}p{(\linewidth - 6\tabcolsep) * \real{0.2500}}
  >{\raggedright\arraybackslash}p{(\linewidth - 6\tabcolsep) * \real{0.2500}}
  >{\raggedright\arraybackslash}p{(\linewidth - 6\tabcolsep) * \real{0.2500}}
  >{\raggedright\arraybackslash}p{(\linewidth - 6\tabcolsep) * \real{0.2500}}@{}}
\toprule\noalign{}
\begin{minipage}[b]{\linewidth}\raggedright
Selected profile
\end{minipage} & \begin{minipage}[b]{\linewidth}\raggedright
\texttt{a\textgreater{}b\textgreater{}c\textgreater{}a}
\end{minipage} & \begin{minipage}[b]{\linewidth}\raggedright
\texttt{a0\textgreater{}b2\textgreater{}c1\textgreater{}a0}
\end{minipage} & \begin{minipage}[b]{\linewidth}\raggedright
Interpretation if a decoder is supplied
\end{minipage} \\
\midrule\noalign{}
\endhead
\bottomrule\noalign{}
\endlastfoot
\nolinkurl{raw_glyphs} & Accepted & Accepted & Character sequence only. \\
\nolinkurl{label_path} & Accepted & Accepted & Whole labels such as
\nolinkurl{a0} are identifiers. \\
\nolinkurl{suffix_path} & Rejected & Accepted & A letter plus a required
digit suffix. \\
\nolinkurl{glyph_path} & Accepted & Rejected & Single glyphs separated by
direction marks. \\
\end{longtable}

An outer enclosure is not a universal requirement. Profile choice is
part of the input contract; a parser must not guess a favorable
interpretation after seeing an otherwise invalid candidate. Byte
sequences, Unicode scalar values, grapheme clusters and grammar tokens
also need separate counting conventions.

The historical inventory describes 22 starts alongside six retained
runtime profiles. That inventory is not evidence of 28 semantic
evaluators or a single shared resource limit. Recognition, AST
construction, decoding, evaluation and exact serialization are distinct
checks. A scalar syntax that accepts \texttt{a:1/0} still needs a
decoder to reject zero denominators unless the grammar itself has been
changed to exclude them.

The source archive preserves the original parser descriptions and
historical codec example. They must be bound to their profile and
implementation revision before being used as current protocol
requirements.

\subsection{11. What a justified lift must
demonstrate}\label{what-a-justified-lift-must-demonstrate}

For a claim that an extension is necessary or minimal:

\begin{enumerate}
\def\labelenumi{\arabic{enumi}.}
\tightlist
\item
  State the failed predicate and protected observations or laws.
\item
  Specify admissible repairs within the current carrier and operations.
\item
  Establish that those repairs cannot meet the predicate, or report
  incomplete search if the budget prevents that conclusion.
\item
  Specify the extension class, equivalence and cost order.
\item
  Retain all tied minima unless an explicit selection policy is applied.
\item
  Check that each accepted extension meets the predicate and preserves
  the particular contract being claimed.
\end{enumerate}

A finite candidate census is exhaustive only for its declared domain. A
budget stop is not proof of obstruction. A declared successor, a
proposed lift and a licensed transition therefore have separate evidence
states. Exact names of runtime terminal values belong to the engineering
protocol's schema; combining papers does not create a new enum or claim
every runner already conforms.

Similarly, proposer, evaluator and authorizer are roles with different
inputs and responsibilities. Three agents are neither necessary nor
sufficient. A theorem that an already halted state is unchanged by
replay establishes that equation; it does not prove the authorization of
an active transition.

\subsection{12. Boundaries retained from the corrected
walkthrough}\label{boundaries-retained-from-the-corrected-walkthrough}

The following withdrawn claims must not re-enter through consolidation:

\begin{longtable}[]{@{}
  >{\raggedright\arraybackslash}p{(\linewidth - 2\tabcolsep) * \real{0.5000}}
  >{\raggedright\arraybackslash}p{(\linewidth - 2\tabcolsep) * \real{0.5000}}@{}}
\toprule\noalign{}
\begin{minipage}[b]{\linewidth}\raggedright
Earlier inference
\end{minipage} & \begin{minipage}[b]{\linewidth}\raggedright
Correct scope
\end{minipage} \\
\midrule\noalign{}
\endhead
\bottomrule\noalign{}
\endlastfoot
Two-sided inverse carrier maps establish univalence. & They establish
the stated equivalence. A univalent universe requires a different
construction and proof. \\
S(n) is automatically a Postnikov tower. & Construction stage and
truncation degree differ; maps and preservation properties of a
fixed-space Postnikov tower are extra data. \\
A selected CHSH expression detects all contextuality. & All required
orientations and the model assumptions must be checked. A Bell score
alone is not a constructed Cech cohomology class. \\
\texttt{F²=I} establishes a pentagon identity. & It is an involution
check; full associator coherence is a different statement. \\
A SQL comparison and a Lean compilation establish bisimulation. & A
relation connecting their states and transitions still needs a proof. \\
Small matrix/graph fixtures certify a physical theory. & They establish
only the checked finite calculation, within their arithmetic model and
input assumptions. \\
\end{longtable}

For the correlator tuple (1,1,−1,1), the selected expression
E00+E01+E10−E11 is zero, while E00+E01−E10+E11 is four. Uniform-marginal
pairwise distributions with those correlations exist. No global
deterministic assignment can satisfy all four signs: the product of the
four required pair products is −1, but every assignment gives +1. This
is the concrete missed orientation from the corrected walkthrough.

The
\href{../../../../../omega_engine_education/documentation/003_minimal_symbol_system_n_mssn/physics_claim_cleanup.md}{physics
cleanup record} records removed fabricated implementations and repaired
finite examples. Valid trace, determinant, commutator, boundary and
finite graph calculations remain useful. Their utility does not
reinstate removed CDT, Super-WDW, sign-problem, chiral-superconductor or
universal physical-verification claims. This paper does not claim a
fresh rerun of those implementations.

\subsection{13. The bridge into WDP}\label{the-bridge-into-wdp}

The common question across the preceding examples is now explicit. Let
q:X→Y be a report and f:X→Z a protected answer. Can the report answer
the question? The criterion is that f be constant on every fiber of q.
Equivalently, there is a unique function g on im(q) with f=g composed
with q.

The image restriction matters. Extending g to all of Y requires values
on unrealized reports; it need not be unique and can fail if the
codomain has no available value. Paper 002 develops report descent,
finite observation costs, linear obstruction, calibrated recovery and
other conditional constructions with their own premises and evidence.

\begin{longtable}[]{@{}
  >{\raggedright\arraybackslash}p{(\linewidth - 2\tabcolsep) * \real{0.5000}}
  >{\raggedright\arraybackslash}p{(\linewidth - 2\tabcolsep) * \real{0.5000}}@{}}
\toprule\noalign{}
\begin{minipage}[b]{\linewidth}\raggedright
Object here
\end{minipage} & \begin{minipage}[b]{\linewidth}\raggedright
Question entering Paper 002
\end{minipage} \\
\midrule\noalign{}
\endhead
\bottomrule\noalign{}
\endlastfoot
A projection of history & Is the protected state constant on its
fibers? \\
A glyph/profile decoder & Which distinctions survive decoding or
normalization? \\
An old-image recovery map & Which source inputs are recoverable, and
where must it abstain? \\
A cochain or quotient representation & Does the protected question
descend, and what witnesses failure? \\
A candidate repair & Does it meet the condition while preserving the
declared observations? \\
\end{longtable}

This is an explanatory bridge with explicit mathematical interfaces. A
complete native derivation of WDP from an opaque inscription remains a
separate claim; this trilogy does not supply it by changing its
publication order.

\subsection{14. Evidence and consolidation
record}\label{evidence-and-consolidation-record}

The \href{paper_claim_ledger.md}{claim ledger} records the scope of each
main result. The \href{source_map.md}{source map} assigns the three
predecessor papers to the sections above and names claims that were
narrowed or omitted. The \href{source_archive/README.md}{source archive}
preserves those manuscripts and their ledgers/change logs. The education
paths now point here; they are not competing active manuscripts.

The \href{source_provenance.json}{source manifest} records the baseline,
original blob identities and digests, and relocated archive digests.
Source links were adjusted for the archive location. Local Lean 4 proofs
live in \href{lean4/Paper001.lean}{lean4}, and the finite Python checks
live in \url{code/sn_mssn_checks.py}. The reproduction entry is
\url{code/reproduce.py}. This revision makes no new model-call,
full-repository coverage, DOI, or publication-deposit claim.

\subsection{Appendix A. The 22 declared MSS grammar
starts}\label{appendix-a.-the-22-declared-mss-grammar-starts}

This table is the source inventory, not a list of completed semantic
adapters. The production grammar remains the authority for accepted
syntax.

\begin{longtable}[]{@{}
  >{\raggedright\arraybackslash}p{(\linewidth - 8\tabcolsep) * \real{0.2000}}
  >{\raggedright\arraybackslash}p{(\linewidth - 8\tabcolsep) * \real{0.2000}}
  >{\raggedright\arraybackslash}p{(\linewidth - 8\tabcolsep) * \real{0.2000}}
  >{\raggedright\arraybackslash}p{(\linewidth - 8\tabcolsep) * \real{0.2000}}
  >{\raggedright\arraybackslash}p{(\linewidth - 8\tabcolsep) * \real{0.2000}}@{}}
\toprule\noalign{}
\begin{minipage}[b]{\linewidth}\raggedright
\#
\end{minipage} & \begin{minipage}[b]{\linewidth}\raggedright
Start Symbol
\end{minipage} & \begin{minipage}[b]{\linewidth}\raggedright
EBNF Production
\end{minipage} & \begin{minipage}[b]{\linewidth}\raggedright
Canonical Input
\end{minipage} & \begin{minipage}[b]{\linewidth}\raggedright
Lexical \& Structural Contract
\end{minipage} \\
\midrule\noalign{}
\endhead
\bottomrule\noalign{}
\endlastfoot
\textbf{01} & \nolinkurl{raw_glyphs} &
\texttt{\{\ glyph\ \textbackslash{}\textbar{}\ "\textless{}"\ \textbackslash{}\textbar{}\ "\textgreater{}"\ \}}
& \texttt{a\textgreater{}\textless{}b0\textless{}\textless{}} & Raw
strings over \(\Gamma_{38}\), including \(\varepsilon\). No balancing
required. \\
\textbf{02} & \nolinkurl{word} & \texttt{glyph,\ \{\ glyph\ \}} &
\nolinkurl{alpha01} & Nonempty sequence over \(\Sigma_{36}\). No
delimiters. \\
\textbf{03} & \nolinkurl{glyph_path} &
\texttt{glyph,\ \{\ direction,\ glyph\ \}} &
\texttt{a\textgreater{}b\textless{}c\textgreater{}a} & Alternating
single glyphs and directions. \\
\textbf{04} & \nolinkurl{label_path} &
\texttt{word,\ \{\ direction,\ word\ \}} &
\texttt{a0\textgreater{}b2\textgreater{}c1} & Directed paths between
multi-glyph identifiers. \\
\textbf{05} & \nolinkurl{suffix_path} &
\texttt{suffix\_vertex,\ \{\ direction,\ suffix\_vertex\ \}} &
\texttt{a0\textgreater{}b12\textless{}c3} & Vertex structured as
\nolinkurl{lower_latin} followed by \nolinkurl{digits}. \\
\textbf{06} & \nolinkurl{path_list} &
\texttt{glyph\_path,\ \{\ ";",\ glyph\_path\ \}} &
\texttt{a\textgreater{}b;c\textless{}d} & Semicolon-delimited sequence
of independent paths. \\
\textbf{07} & \nolinkurl{suffix_list} &
\texttt{suffix\_path,\ \{\ ";",\ suffix\_path\ \}} &
\texttt{a0\textgreater{}b1;c2\textless{}d3} & Semicolon-delimited
sequence of suffix paths. \\
\textbf{08} & \nolinkurl{scalar_path} &
\texttt{scalar\_vertex,\ \{\ direction,\ scalar\_vertex\ \}} &
\texttt{a:1.5\textgreater{}b:-2/3} & Paths with colon-separated scalar
weights. \texttt{a:1/0} parses syntactically. \\
\textbf{09} & \nolinkurl{vector_path} &
\texttt{vector\_vertex,\ \{\ direction,\ vector\_vertex\ \}} &
\texttt{a{[}1,2{]}\textgreater{}b{[}3,4{]}} & Bracketed vector-annotated
vertices. \\
\textbf{10} & \nolinkurl{edge_weight_path} &
\texttt{lower\_latin,\ \{\ weighted\_link,\ lower\_latin\ \}} &
\texttt{a\textgreater{}2.5\textgreater{}b\textless{}1\textless{}c} &
Arrow-flanked edge weights (\texttt{\textgreater{}w\textgreater{}} or
\texttt{\textless{}w\textless{}}). Opposing marks
(\texttt{\textgreater{}w\textless{}}) rejected. \\
\textbf{11} & \nolinkurl{labelled_path} &
\texttt{lower\_latin,\ \{\ "-",\ word,\ "\textgreater{}",\ lower\_latin\ \}}
& \texttt{a-step\textgreater{}b} & Named relation links
(\texttt{-label\textgreater{}}). \\
\textbf{12} & \nolinkurl{prefix_tree} &
\texttt{glyph\ \textbackslash{}\textbar{}\ direction,\ prefix\_tree,\ prefix\_tree}
& \texttt{\textgreater{}ab} & Polish prefix notation for binary syntax
trees. \\
\textbf{13} & \nolinkurl{postfix_tree} &
\texttt{glyph,\ \{\ postfix\_tree,\ direction\ \}} &
\texttt{ab\textgreater{}} & Reverse Polish notation for binary syntax
trees. \\
\textbf{14} & \nolinkurl{sequence} & \texttt{glyph,\ \{\ ",",\ glyph\ \}} &
\texttt{a,b,c} & Comma-delimited list of single glyphs. \\
\textbf{15} & \nolinkurl{nested_tree} &
\texttt{glyph\ \textbackslash{}\textbar{}\ "(",\ {[}\ nested\_tree,\ \{\ ",",\ nested\_tree\ \}\ {]},\ ")"}
& \texttt{(a,(b,c))} & Parenthesized nested trees. \\
\textbf{16} & \nolinkurl{head_term} &
\texttt{word,\ {[}\ "(",\ {[}\ head\_term,\ \{\ ",",\ head\_term\ \}\ {]},\ ")"\ {]}}
& \texttt{f(x,g(y))} & Functional terms with head words and argument
lists. \\
\textbf{17} & \nolinkurl{incidence} &
\texttt{incidence\_item,\ \{\ ";",\ incidence\_item\ \}} &
\texttt{{[}a,b{]}\textgreater{}{[}c,d{]}} & Hypergraph directed
incidence lists. \\
\textbf{18} & \nolinkurl{unicode_path} &
\texttt{unicode\_glyph,\ \{\ unicode\_direction,\ unicode\_glyph\ \}} &
\texttt{α→β←γ} & Directed paths over Greek/symbolic Unicode
alphabets. \\
\textbf{19} & \nolinkurl{arrow_path} &
\texttt{glyph,\ \{\ unicode\_direction,\ glyph\ \}} & \texttt{a→b←c} &
ASCII glyphs joined by Unicode arrows (\texttt{→}, \texttt{←}). \\
\textbf{20} & \nolinkurl{bitstream} & \texttt{bit,\ \{\ bit\ \}} &
\nolinkurl{10011010} & Nonempty binary digit strings (\nolinkurl{0} and
\nolinkurl{1}). \\
\textbf{21} & \nolinkurl{octets} &
\texttt{hex\_digit,\ hex\_digit,\ \{\ hex\_digit,\ hex\_digit\ \}} &
\nolinkurl{deadbeef} & Lowercase hexadecimal byte strings (even length
required). \\
\textbf{22} & \nolinkurl{rewrite_rule} &
\texttt{\{\ glyph\ \},\ "=\textgreater{}",\ \{\ glyph\ \}} &
\texttt{a=\textgreater{}b} & String rewrite rules with optional empty
sides. \\
\end{longtable}

\end{document}
